\subsection{Holomorphic Functions}

In this section, we assume that K = C.

3.3.1. Suppose F is quasi-complete. Let U be an open subset of E and f a mapping from U into F. The following conditions are equivalent:

\begin{enumerate}
\def\labelenumi{(\roman{enumi})}
\item
  \(f\) is holomorphic;
\item
  \(f\) is differentiable;
\item
  \(f\) is locally bounded and for any \(a \in \mathbf{U}\) and \(h \in \mathbf{E}\) , the function \(t \mapsto f(a + th)\) defined in an open neighborhood of 0 in \(\mathbf{C}\) is holomorphic;
\item
  \(f\) is locally bounded and for every continuous linear form \(u\) on F, the function \(u \circ f\) with values in C is holomorphic;
\item
  \(f\) is continuous and locally bounded, and there exists a total set H in the dual of F such that \(u \circ f\) is holomorphic for every \(u \in \mathbf{H}\) .
\end{enumerate}

When E is finite-dimensional (resp. when F is a Banach space), these conditions are still equivalent to conditions (iii′), (iv′), or (v′) (resp. (iv′) or (v′)) obtained from (iii), (iv), or (v) (resp. (iv) or (v)) by removing the hypothesis “f is locally bounded.”

3.3.2. Suppose F is quasi-complete. Let U be an open subset of E and \((f_{n})\) a sequence of holomorphic maps from U into F, having the following property:

\begin{enumerate}
\def\labelenumi{(\Alph{enumi})}
\setcounter{enumi}{22}
\tightlist
\item
  Every point of U has a neighborhood in which the sequence \((f_{n})\) converges uniformly.
\end{enumerate}

Then the limit \(f\) of the sequence \((f_{n})\) is holomorphic, the sequence of derivatives \((\mathrm{D}f_{n})\) (with values in the quasi-complete space \(\mathcal{L}(\mathrm{E};\mathrm{F})\) ) possesses property (W) and Df is the limit of \((\mathrm{D}f_{n})\) .

3.3.3. Let U be an open subset of E and f a holomorphic map from U into F, with F assumed to be quasi-complete. Let \(\mathbf{R} = (\mathbf{R}_{i}) \in (\mathbf{R}_{+}^{*})^{n}\) and suppose that the polyball \(\mathbf{B}(\mathbf{R})\) is contained in U and f is bounded on \(\mathbf{B}(\mathbf{R})\) . We then have, for every \(\alpha \in N^{n}\) and every \(x = (x_{i}) \in \mathbf{B}(\mathbf{R})\) :

\[
\Delta^ {\alpha} f (0) (x) = \int_ {0} ^ {1} \dots \int_ {0} ^ {1} f (\mathbf {e} (\theta_ {1}) x _ {1}, \dots , \mathbf {e} (\theta_ {n}) x _ {n}) \mathbf {e} (- \alpha_ {1} \theta_ {1} - \dots - \alpha_ {n} \theta_ {n}) d \theta_ {1} \dots d \theta_ {n}
\]

\((\text{with } \mathbf{e}(\theta) = \exp 2\pi i \theta)\) .

Additionally, let \(\gamma\) be a continuous seminorm on \(F\) and let \(M\) be the supremum of \(\|f(x)\|_{\gamma}\) for \(\|x_{i}\|=R_{i}\) . One has \(\|\Delta^{\alpha}f(0)(x)\|_{\gamma}\leqslant M\) for all \(x\in B(R)\) and \(\|\Delta^{\alpha}f(0)\|_{\gamma}\leqslant MR^{-\alpha}\) . Finally, the domain of convergence of the series expansion of \(f\) at 0 contains the interior of the polydisc \(B(R)\) .

3.3.4. We keep the hypotheses of 3.3.3 and suppose moreover that \(\mathbf{E}_i = \mathbf{C}\) . Let \(\sum X^{\alpha}c_{\alpha}\) be the series expansion of \(f\) at 0. We have:

\[
c _ {\alpha} = \mathbf {R} ^ {- \alpha} \int_ {0} ^ {1} \dots \int_ {0} ^ {1} f (\mathbf {e} (\theta_ {1}) \mathrm{R} _ {1}, \dots , \mathbf {e} (\theta_ {n}) \mathrm{R} _ {n}) \mathbf {e} (- \alpha_ {1} \theta_ {1} - \dots - \alpha_ {n} \theta_ {n}) d \theta_ {1} \dots d \theta_ {n}
\]

and:

\[
\| c _ {\alpha} \| _ {\gamma} \leqslant \mathbf {R} ^ {- \alpha} \sup _ {x \in \mathbf {B} (\mathbf {R})} \| f (x) \| _ {\gamma}
\]

("Cauchy inequalities"). The strict domain of convergence of the series \(\sum_{\alpha} X^{\alpha} c_{\alpha}\) contains the interior of B(R).

3.3.5. Suppose E is finite-dimensional and F is quasi-complete. There then exists in \(\mathcal{H}(\mathbf{E};\mathbf{F})\) a series \(f_{0}\), one and only one, with infinite radius of convergence (for every norm on E), such that \(f(x)=f_{0}(x)\) for all \(x\in E\) .

3.3.6. If \(f\) is a holomorphic mapping from E into F such that \(f(E)\) is bounded, then the function \(f\) is constant (“Liouville’s theorem”).

3.3.7. We assume that \(E \neq 0\) . Let f be a holomorphic mapping from an open set U of E into F. Let \(a\) be a point of U and \(\gamma\) be a continuous seminorm on F. For every neighborhood V of a, contained in U, there exists \(x \in V\) , \(x \neq a\) , such that:

\[
\| f (a) \| _ {\gamma} \leqslant \| f (x) \| _ {\gamma}.
\]

If moreover F = C and if f is not constant in a neighborhood of a, we have \(|f(a)| < \sup_{x \in V, x \neq a} |f(x)|\) and the map f is open in a neighborhood of a. Finally, let B be a bounded open set whose closure is contained in U, and let B' be its

boundary. We have:

\[
\sup _ {x \in \overline {{\mathbf {B}}}} | f (x) | = \sup _ {x \in \mathbf {B} ^ {\prime}} | f (x) |
\]

(“maximum principle”).

3.3.8. Suppose E is finite-dimensional. Let U be an open subset of E, and let S be a vector subspace of codimension \(\geqslant2\) of E and f a holomorphic mapping from \(\mathrm{U}\cap(\mathrm{E}-\mathrm{S})\) into F. Then f extends continuously to a holomorphic mapping from U into F.

Suppose that E = C and let \(0 \leqslant \lambda < \mu \leqslant +\infty\) . Let f be a holomorphic mapping from the open set \(U = \{z \in C | \lambda < |z| < \mu\}\) into F. There exists a unique family \((a_{n}(f))_{n \in \mathbb{Z}}\) of elements of F such that, for every compact H of U, the family \((a_{n}(f)z^{n})_{n \in \mathbb{Z}}\) is uniformly summable, with sum \(f(z)\) as z ranges over H ("Laurent expansion").

Suppose further that \(\lambda=0\) . We call the order of \(f\) at the point \(x=0\) the lower bound (finite or infinite) of the set of integers \(n\) such that \(a_{n}(f)\neq0\) . If there exists a neighborhood V of 0 such that \(f\) is bounded in V- \(\{0\}\) , then \(f\) is of order 0 at the point \(x=0\) and extends by continuity to a holomorphic function on the open set \(|z|<\mu\) . Let \(p\) be an integer \(>0\) ; if \(f\) is of order \(-p\) at the point \(x=0\) , one says that 0 is a pole of order \(p\) of \(f\) .

3.3.10. Suppose that E = C and that F is normed. Let f be a holomorphic mapping from the open unit disk U of E into F, such that \(f(0) = 0\) and let \(M = \sup_{z \in U} \|f(z)\|\) . One then has \(\|f(z)\| \leqslant M \cdot |z|\) for all \(z \in U\) (“Schwarz lemma”).

